\section{Computation and implementation}\label{sec:computation}

Throughout this section $n$ is the number of records, $p$ the number of columns of the model matrix
$X$, intercept included, $G$ the number of groups and $d$ the number of basis columns. In the designs of
this paper $p$ and $d$ are small: $d$ is $1$ or $3$, or one more in the external mode.

\subsection{The algorithm and its cost}\label{sec:algorithm}

We take as given the fitted logistic model, that is, the fitted risks $\hat\pi_i$, the model
matrix $X$ and the fitted weights $W$. Every test of the fitted model needs this fit, and EDGE needs it
only once. Given the fit, EDGE is computed in seven steps.

\begin{enumerate}
\item \textbf{Sort and group.} Rank the records by fitted risk, ties in the order of the records, and
assign the record of rank $R_i$ to group $\lceil R_iG/n\rceil$. A stable sort costs $O(n\log n)$
operations. This is the only step whose cost grows faster than linearly in $n$.
\item \textbf{Group sums and grouped residuals.} In one pass over the records, accumulate for each
group the observed events $O_g$, the expected events $E_g$, the binomial variance $V_g$ and the count
of records, and from them $\pibar_g$ and the residual $r_g$ of \eqref{eq:resid}, at a cost of $O(n)$.
\item \textbf{Estimation adjustment.} In the same pass, accumulate the $G\times p$ matrix $U$ of
\eqref{eq:omega}, at a cost of $O(np)$. Form the information matrix $X^\top WX$, at a cost of $O(np^2)$,
and factor it, at a cost of $O(p^3)$. The $G\times G$ covariance $\Omega$ is not formed, because the
test needs only its products with the $d$ basis columns (step~6). In the external mode this step is
skipped, because $\Omega=I_G$.
\item \textbf{Basis.} Evaluate the $d$ shape functions of Section~\ref{sec:bases} at the $G$ group
means $\pibar_g$, and orthogonalise the columns of $Z$, at a cost of $O(Gd^2)$.
\item \textbf{Projection.} Compute $Z^\top r$, at a cost of $O(Gd)$, and the statistic $S$ of
\eqref{eq:stat} by solving a $d\times d$ system, at a cost of $O(d^3)$.
\item \textbf{Eigenvalues.} Compute the $d\times d$ matrix $(Z^\top Z)^{-1}Z^\top\Omega Z$ of
Theorem~\ref{thm:null} as $Z^\top\Omega Z=Z^\top Z-(Z^\top U)(X^\top WX)^{-1}(U^\top Z)$, which costs
$O(Gdp+dp^2+d^2p)$ and never forms $\Omega$, and its eigenvalues, at a cost of $O(d^3)$. Forming $\Omega$
first would cost $O(G^2(p+d))$ operations and $O(G^2)$ memory.
\item \textbf{Reference distribution.} Match the first two moments of the weighted sum in
\eqref{eq:wchi2} by a scaled chi-squared variable $c\,\chi^2_\nu$, with
$\nu=(\sum_j\lambda_j)^2/\sum_j\lambda_j^2$ and $c=\sum_j\lambda_j^2/\sum_j\lambda_j$, and report
$p=\Pr(\chi^2_\nu>S/c)$ \citep{satterthwaite1946}. This costs $O(d)$ operations and one evaluation of
the chi-squared distribution function.
\end{enumerate}

The total cost, given the fit, is therefore
\[
O\bigl(n\log n+np^2+p^3+Gd(d+p)+dp^2+d^3\bigr)=O\bigl(n\log n+np^2\bigr),
\]
because $G\le n$, $p\le n$ and $d$ is bounded. The sort dominates only when $p^2$ is small next to
$\log n$; otherwise forming the information matrix dominates, as it does in the fit itself. The memory
needed beyond the data is $O(Gp)$.

The avoided $G\times G$ term matters at the default partition, $G=\max(10,\lceil n/25\rceil)$, where it
grows with the square of $n$; at $n=100{,}000$ there are $4000$ groups. The timed implementation of
EDGE forms no $G\times G$ matrix and is in the reproducibility archive; the CRAN release of the package
forms $\Omega$. The two return identical results up to rounding: over $250$ data sets and $6600$ settings, covering
the three bases, both weightings, both partitions and both modes, the largest differences were
$9.3\times10^{-10}$ in the statistic and $1.4\times10^{-11}$ in the $p$-value. Forming $\Omega$ made
EDGE $1.6$ times slower at $n=500$ and $11$ times slower at $n=100{,}000$, where it took $0.47$ seconds
against $43$ milliseconds.

\subsection{A reference distribution without resampling or refitting}\label{sec:noresampling}

The weights of the reference distribution of Theorem~\ref{thm:null} are the eigenvalues of a $d\times
d$ matrix built from quantities that the fit has already produced. No second model fit, no bootstrap
sample and no random split is needed. Three consequences follow.

First, the test costs about as much as the other tests that need no resampling. Every timing was taken on one
core of an AMD Ryzen 9 3900X processor with $32$ gigabytes of memory, in R 4.4, with the linear-algebra
library held to one thread and nothing else running. At $n=1000$, EDGE takes $2$ milliseconds, Stukel's joint
score test $1$, the Hosmer--Lemeshow statistic $2$, the cubic calibration test $5$ and the GiViTI belt
$9$ (Figure~\ref{fig:cost}). At $n=100{,}000$ EDGE takes $43$ milliseconds, Stukel's score $18$ and the
Hosmer--Lemeshow statistic $120$. The Liu projection test draws $B=250$ bootstrap samples and takes
$35$ seconds at $n=1000$. BAGofT refits the model $(1+\textit{nsim})\times\textit{nsplits}$ times,
$10{,}100$ model fits at its defaults; it takes $392$ seconds at $n=1000$ and ten minutes at $n=2000$.
Le Cessie's smoothed test needs no resampling, but in its published form it multiplies $n\times n$
matrices, so its cost grows with the cube of $n$: $5.6$ seconds at $n=1000$, $74$ seconds at $n=2000$
and $21$ minutes at $n=5000$, where EDGE takes $4$ milliseconds. Its $n\times n$ kernel would need about
$10$ gigabytes on the $35{,}000$ patients of each half of the cohort of Section~\ref{sec:realdata}, so
it was not run there.

Second, the $p$-value is a deterministic function of the data. It carries no Monte Carlo error and
does not depend on a random seed, whereas the two resampling tests return $p$-values on grids of
$1/100$ and $1/250$.

Third, the test can be run wherever it is needed: on every replicate of every scenario of this paper,
where the three slow rivals were run only on subsets, and across partitions. Because its degrees of
freedom do not depend on $G$ (Theorem~\ref{thm:growingG}), a sweep over partitions costs one sort in
all and one pass for each partition (Section~\ref{sec:realpartition}).

\subsection{The external mode and the streaming form}\label{sec:streamcomp}

In the external mode $\Omega=I_G$ and a constant column is added to the basis
(Section~\ref{sec:null}). Step~3 of the algorithm disappears, and the model matrix is not needed. All
$d+1$ eigenvalues of step~6 equal one, so the Satterthwaite approximation of step~7 returns the
chi-squared distribution on $d+1$ degrees of freedom exactly. The cost is $O(n\log n+Gd^2)$. On one
core, the external mode at ten groups took about $0.3$ milliseconds at $n=1000$ and $12$ milliseconds at
$n=100{,}000$. The test of \citet{sadatsafavi2024nonparametric}, from the R package \texttt{cumulcalib},
took about $0.25$ and $14$ milliseconds on the same data sets.

The external mode can also be computed as a stream, which suits the monitoring of a deployed model.
The equal-frequency groups of step~1 are replaced by cut points $0<c_1<\dots<c_{G-1}<1$ fixed in
advance from reference predictions, without any outcome. The monitor holds four sums for each group:
the observed events $O_g$, the expected events $E_g$, the binomial variance $V_g$ and the count $m_g$.
These are all the statistic needs, because $r_g=(O_g-E_g)/\sqrt{V_g}$ and $\pibar_g=E_g/m_g$.

\begin{enumerate}
\item \textbf{Update.} When a patient with prediction $p_i$ and outcome $y_i$ arrives, find the
patient's group by a binary search among the $G-1$ cut points, at a cost of $O(\log G)$, and add $y_i$,
$p_i$, $p_i(1-p_i)$ and $1$ to the four sums of that group. No earlier record is kept or revisited.
\item \textbf{Test.} At any time, compute the residuals and the group means from the $G$ summaries,
evaluate the basis with its constant column over the non-empty groups, and compute $S_n=r^\top P_Zr$
through a QR decomposition of a $G\times k$ matrix, where $k$ is the rank of the basis ($d+1$ once
$d+1$ groups with distinct mean risks are non-empty), at a cost of $O(Gk^2)$.
\end{enumerate}

The memory needed is $O(G)$, whatever the length of the stream. Each of the four sums is a sum of terms
that depend on one patient alone, so the streamed statistic is the same for every order and every
division into batches, and equals the external-mode statistic computed in one pass with the same cut
points. Under a calibrated model the statistic has a chi-squared reference on $k$ degrees of freedom
asymptotically, as every non-empty group grows, however unequal the groups become (Online Resource~1,
Section~S8, proves these three properties). The price of fixed cut points is that the groups drift
away from equal size when the risk mix of the monitored patients changes, and small groups give weaker
protection against corrupted records (Proposition~\ref{prop:bounded}). On one core, after a million
patients had been added, adding a hundred more took about half a millisecond and computing the test
about $1.5$ milliseconds.

In a simulation of the streaming form fixed before it was run, at ten groups and with $2000$
replicates, the level was $0.047$ and $0.054$ after $10{,}000$ and $100{,}000$ patients, and $0.050$ and
$0.052$ when every monitored patient came from a risk distribution different from the reference. The
test detected an overfitted model, calibration slope $0.8$, in $0.99$ of replicates after $2000$
patients (Online Resource~1, Section~S8). The statistic accumulates since deployment; to follow drift,
the summaries are restarted each period and each period is tested on its own. One test at any time is
valid, but acting on the first rejection over repeated looks is not: ten looks at $5\%$ raised the
false-alarm rate to $0.21$, and splitting the level over ten looks fixed in advance, $0.5\%$ each,
brought it to $0.030$.

\subsection{Software}\label{sec:software}

The test and its comparators are implemented in the R package \texttt{ebrahim.gof}
\citep{ebrahim2025gof}, available from CRAN (Section~\ref{sec:algorithm} compares its timing with the
implementation timed here).

\begin{itemize}
\item \texttt{edge.gof(fit)} runs EDGE on a fitted binomial \texttt{glm} and reports the verdict at the
default partition and, as a robustness check, at ten groups; the degrees of freedom reported are the
effective $\nu$ of step~7. The argument \texttt{G = 10} gives the ten-group partition, \texttt{basis =
"sym"} or \texttt{"stukel"} the other bases, and \texttt{weights = "score"} the score form.
\item \texttt{edge.gof(y, predicted\_probs = p, external = TRUE, G = 10)} runs the external mode on
outcomes \texttt{y} and frozen predictions \texttt{p}; with \texttt{G} set accordingly it reproduces
every statistic of Section~\ref{sec:realdata}.
\item \texttt{run.all.gof()} fits the model once and returns every comparator of
Section~\ref{sec:comparators} except the cubic calibration test.
\item \texttt{run.all.external(y, p)} runs the external mode beside the other external-validation
tests, Cox's recalibration test and the calibration intercept and slope among them.
\item \texttt{edge.stream()} creates the monitor of Section~\ref{sec:streamcomp}; \texttt{update()} adds
a batch of outcomes and predictions and \texttt{summary()} returns the test at any time.
\end{itemize}
